\documentclass[11pt]{article}
\usepackage[a4paper,margin=1in]{geometry}
\usepackage{fontspec}
\setmainfont{DejaVu Serif}
\usepackage{microtype}
\usepackage{amsmath,amssymb,mathtools}
\setlength{\parindent}{1.2em}
\setlength{\parskip}{0.25em}
\emergencystretch=3em
\sloppy
\begin{document}

\begin{center}
{\Large G. Lejeune Dirichlet's Werke. II.}\par\vspace{1.2em}
{\Large XXIV.}\par\vspace{0.5em}
{\LARGE \textbf{Ueber den biquadratischen Charakter der Zahl „Zwei“}}\par\vspace{1em}
{\large Von G. Lejeune Dirichlet}\par\vspace{0.5em}
{\small Aus einem Briefe Dirichlet's an Herrn Stern zu Göttingen, Crelle's Journal für die reine und angewandte Mathematik, Bd. 57 S. 187--188.}
\end{center}

\bigskip
\begin{center}
{\large \textbf{Ueber den biquadratischen Charakter der Zahl „Zwei“.}}
\end{center}

Es sei $p$ eine Primzahl der Form $4n+1$ und
\[
p=a^2+b^2,
\]
wo $a$ ungerade. Aus dieser Gleichung und der daraus folgenden
\[
2p=(a+b)^2+(a-b)^2
\]
ergiebt sich sogleich mit Anwendung des verallgemeinerten Symbols von Legendre
\[
\left(\frac{p}{a}\right)=1, \qquad \left(\frac{2p}{a+b}\right)=1
\]
oder
\[
\left(\frac{p}{a+b}\right)=\left(\frac{2}{a+b}\right),
\]
und dann nach bekannten Sätzen:
\[
\left(\frac{a}{p}\right)=1, \qquad \left(\frac{a+b}{p}\right)=(-1)^{\frac18[(a+b)^2-1]}
\]
oder, was dasselbe ist:
\[
a^{\frac12(p-1)}\equiv 1, \qquad (a+b)^{\frac12(p-1)}\equiv (-1)^{\frac18[(a+b)^2-1]} \quad (\mathrm{mod}.\,p).
\]
Andererseits erhält man aus:
\[
(a+b)^2 \equiv a^2+b^2+2ab \equiv 2ab \quad (\mathrm{mod}.\,p)
\]
durch Erhebung zur Potenz $\frac14(p-1)$:
\[
(a+b)^{\frac12(p-1)}\equiv 2^{\frac14(p-1)}a^{\frac14(p-1)}b^{\frac14(p-1)} \quad (\mathrm{mod}.\,p)
\]
oder, wenn man
\[
b\equiv af
\]
setzt und die beiden letzten Congruenzen berücksichtigt:
\[
(-1)^{\frac18[(a+b)^2-1]}=(-1)^{\frac18(p-1+2ab)}\equiv 2^{\frac14(p-1)}f^{\frac14(p-1)} \quad (\mathrm{mod}.\,p).
\]
Nun ist, weil
\[
a^2+b^2=p, \qquad b^2\equiv a^2f^2,
\]
also
\[
f^2\equiv -1 \quad (\mathrm{mod}.\,p),
\]
\[
(f^2)^{\frac18(p-1+2ab)}=f^{\frac14(p-1)}\cdot f^{\frac12 ab}\equiv 2^{\frac14(p-1)}\cdot f^{\frac14(p-1)} \quad (\mathrm{mod}.\,p)
\]
oder, wenn man mit $f^{\frac14(p-1)}$ dividirt,
\[
2^{\frac14(p-1)}\equiv f^{\frac12 ab} \quad (\mathrm{mod}.\,p),
\]
welches Resultat in das von Gauss im § 24 seiner „theoria residuorum biquadraticorum“ gegebene übergeht, wenn man mit $a^{\frac12ab}$ multiplicirt und für $af$ wieder $b$ einführt.

\begin{center}
Göttingen, den 21. Januar 1857.
\end{center}

\end{document}
